\documentclass[11pt]{article}
\usepackage[margin=24mm]{geometry}
\usepackage{amsmath,amssymb,amsthm,mathtools}
\usepackage[hidelinks]{hyperref}
\newcommand{\tr}{\operatorname{tr}}
\newcommand{\Rea}{\operatorname{Re}}
\title{The same finite heat action: a Dyson remainder and a rational stopping rule}
\author{}\date{20 September 2026}
\begin{document}\maketitle

The source of the finite trace-action question is Connes and van
Suijlekom \cite[Section on the spectral action]{CV}; the original
arithmetic Gaussian context is Connes \cite[Gaussian equation ft]{CH}.
Every finite estimate used here is proved below. An incoming independent
research continuation, Track III, H4--H16, gives these same Dyson and
rational-tail formulas. This note records their complete proof and their
exact agreement with the accompanying NH1--36 and HM1--19 proofs,
without counting the overlapping heat-gap theorem twice.
These local derivations are not substitutes for the cited human sources,
and no infinite arithmetic spectral interpretation is assumed.

\section{Original coordinates and the two observations}
Fix the attained arithmetic quotient with its specified positive
Hermitian matrix \(G\), adjoint \(X^\dagger=G^{-1}X^*G\), the original
selfadjoint compression \(A=A^\dagger\), and its rank-one relation
\[
 Q=r\ell,\qquad \ell r=0,\qquad
 \epsilon=\|Q\|_G=\|Q\|_{\mathrm{HS},G},\qquad T(t)=A-tQ .
 \tag{DS1}
\]
Here \(t\) is the outgoing-relation parameter, not de Bruijn--Newman
time. For real \(t\) and positive heat time \(s\), retain
\[
 \Phi_s(t)=\tr e^{-sT(t)^2},\quad
 \Psi_s(t)=\tr e^{-sT(t)^\dagger T(t)},\quad
 \Delta_s(t)=\Rea\Phi_s(t)-\Psi_s(t).
 \tag{DS2}
\]
At the quotient endpoint \(t=1\), \(T(1)=M=(M_S-cI)/i\), where
\(c=k/2\). Therefore the holomorphic heat test in the original scaling
coordinate is \(e^{s(M_S-cI)^2}\); its singular-value counterpart is
\(e^{-s(M_S-cI)^\dagger(M_S-cI)}\). No centre or factor \(i\) is dropped.
The original physical unit and all root multiplicities are those in
NH1--6 and NH24--26.

\section{A second convergent expansion of the same action}
Since \(Q^2=0\), put \(D=AQ+QA\). Then
\[
 T(t)^2=A^2-tD,\qquad
 \|D\|_G\le 2\|A\|_G\epsilon,\qquad
 \|D\|_{1,G}\le 2\|A\|_G\epsilon .
 \tag{DS3}
\]
The last bounds follow by the triangle inequality and the rank-one
trace-norm identities \(\|AQ\|_{1,G}\le\|A\|_G\|Q\|_{1,G}\) and
\(\|Q\|_{1,G}=\epsilon\). In particular no factor equal to the quotient
dimension enters.
For the simplex
\(\Delta_n=\{(v_0,\ldots,v_n):v_j\ge0,\sum v_j=1\}\), use its usual
ordered-time measure, whose volume is \(1/n!\). Then
\[
 \Phi_s(t)=\tr e^{-sA^2}+\sum_{n\ge1}t^n d_n(s),\qquad
 d_n(s)=s^n\int_{\Delta_n}
 \tr\!\left(e^{-sv_0A^2}D e^{-sv_1A^2}\cdots D e^{-sv_nA^2}\right)dv .
 \tag{DS4}
\]
This is an entire series in complex \(t\). To prove it, the matrix
exponential solves
\[
 U(s)=e^{-sA^2}
       +t\int_0^s e^{-(s-v)A^2}D U(v)\,dv.
 \tag{DS5}
\]
Iterating this Volterra equation yields precisely the ordered simplex
integrals in DS4, after \(v\mapsto sv\). Selfadjointness gives
\(\|e^{-vA^2}\|_G\le1\) for \(v\ge0\).
Taking one factor \(D\) in trace norm and all remaining factors in
operator norm proves
\[
 |d_n(s)|\le \frac{s^n\|D\|_{1,G}\|D\|_G^{n-1}}{n!}.
 \tag{DS6}
\]
The corresponding operator-norm estimates prove uniform absolute
convergence on every bounded complex \(t\)-disk. The convergent series
solves DS5 and hence equals the exponential by uniqueness of the finite
matrix differential equation. If \(D=0\), all the terms vanish.
Otherwise, for every integer \(m\ge0\), with \(x=s|t|\|D\|_G\),
\[
 \left|\Phi_s(t)-\tr e^{-sA^2}-\sum_{n=1}^m t^nd_n(s)\right|
 \le s|t|\|D\|_{1,G}\frac{e^x x^m}{(m+1)!}.
 \tag{DS7}
\]
Indeed \(\sum_{n\ge m+1}x^{n-1}/n!
\le x^m e^x/(m+1)!\), by factoring the first term and bounding each
subsequent factorial ratio. This expansion and NH7--16 are Taylor
series of the same entire finite function at \(t=0\), so their
coefficients agree. In particular the cyclic divided-difference
coefficient remains \((-1)^n/n\), with all confluent nodes and complex
residues retained as in the complete proof NH7--10; it is not \(n!\).

\section{Full proof of a rational heat-gap remainder}
For any finite \(T\), put \(a_n=\tr((T^\dagger T)^n)-\Rea\tr T^{2n}\).
Choose a \emph{specified upper bound} \(B\ge\|T\|_G^2\).
For example \(B=\tr(T^\dagger T)\) is always valid. It is not asserted
that an arbitrary chosen upper bound is less than that trace.
We have
\[
 a_1=\tfrac12\|T-T^\dagger\|_{\mathrm{HS},G}^2,\qquad
 |a_n|\le n(2n-1)B^{n-1}a_1 .
 \tag{DS8}
\]
For completeness write \(T=H+iK\) with selfadjoint \(H,K\), and
\(T_v=H+ivK=(1+v)T/2+(1-v)T^\dagger/2\), \(0\le v\le1\).
Its norm is at most \(\sqrt B\).
The two traces defining
\(f_n(v)=\tr((T_v^\dagger T_v)^n)-\Rea\tr T_v^{2n}\)
agree at zero and are even in \(v\), by adjunction and cyclicity.
Thus \(f_n(0)=f_n'(0)=0\).
Each second derivative has \(2n(2n-1)\) ordered choices of two
differentiated factors. Each resulting trace contains two \(K\)'s,
and is bounded by
\(\|K\|_{\mathrm{HS},G}^2B^{n-1}\), using cyclicity and
Hilbert--Schmidt Cauchy--Schwarz. There are two original traces, so
\[
 |f_n''(v)|\le4n(2n-1)\|K\|_{\mathrm{HS},G}^2 B^{n-1}.
\]
Integrating against \(1-v\) gives DS8, since
\(a_1=2\|K\|_{\mathrm{HS},G}^2\).
This also proves the coefficient estimate used in HM9--14.
Expanding the two entire exponentials now gives
\[
 \Delta_s(T)=\sum_{n\ge1}\frac{(-1)^{n+1}s^na_n}{n!},
 \qquad
 |\Delta_s(T)-sa_1|\le sa_1\{(1+2x)e^x-1\},\quad x=sB.
 \tag{DS9}
\]
For an integer \(m\ge1\) with \(m+1>3x\), a completely rational bound is
\[
 \left|\Delta_s(T)-\sum_{n=1}^m
             \frac{(-1)^{n+1}s^na_n}{n!}\right|
 \le sa_1\frac{(2m+1)x^m}{m!\{1-3x/(m+1)\}} .
 \tag{DS10}
\]
To verify every index, the omitted majorant begins at
\(b_{m+1}=(2m+1)x^m/m!\), where
\(b_n=(2n-1)x^{n-1}/(n-1)!\).
For \(n\ge m+1\),
\[
 \frac{b_{n+1}}{b_n}
 =\frac{x(2n+1)}{n(2n-1)}
 \le\frac{3x}{n}\le\frac{3x}{m+1}<1.
\]
Summing this geometric majorant proves DS10. The zero cases \(x=0\)
or \(a_1=0\) follow directly from DS8, without division by them.
When \(T,G,s,B\) have rational complex or rational entries as appropriate,
DS10 supplies a rational interval with no sampled exponential.

For the original path DS1 and real \(t\), \(a_1=t^2\epsilon^2\):
expand \(\|Q-Q^\dagger\|_{\mathrm{HS},G}^2\), use
\(\tr Q^2=\tr(Q^\dagger)^2=0\), and obtain \(2\epsilon^2\).
At \(x\le1/8\), the elementary positive-series inequality
\(e^x\le(1-x)^{-1}\le8/7\) in DS9 gives
\[
 \frac47 st^2\epsilon^2\le\Delta_s(t)
 \le\frac{10}{7}st^2\epsilon^2 .
 \tag{DS11}
\]
This and NH35's half-to-three-halves comparison on \(x\le1/7\)
are two explicit consequences of the same bound DS9, not distinct
independent arithmetic estimates.

\section{Exact metric transport and its finite scope}
For an invertible original conductor \(L:U\to W\), keep
\(T_W=LT_UL^{-1}\) and both specified metrics. Define
\(H=L^*G_WL\). Direct substitution in the adjoint gives
\[
 T_W^{\dagger_{G_W}}=L T_U^{\dagger_H}L^{-1},\qquad
 \tr e^{-sT_W^{\dagger_{G_W}}T_W}
 =\tr e^{-sT_U^{\dagger_H}T_U}.
 \tag{DS12}
\]
This is the metric carried by the positive heat observation; the
holomorphic heat trace is invariant under similarity without choosing
a new metric. No substitution \(H=G_U\) is made. The complete
nonidentity-metric examples and recoverable selfadjoint dilation appear
in NH24--26 and HM15--19. All norms above are evaluated in the metric
actually written. These finite formulas do not bound the growth of the
original arithmetic allowance as \(k\) or the polynomial cutoff grows.

\begin{thebibliography}{9}
\bibitem{CV} Alain Connes and Walter D. van Suijlekom,
\emph{Quadratic Forms, Real Zeros and Echoes of the Spectral Action},
arXiv:2511.23257v1 (2025), Section on the spectral action.
\url{https://arxiv.org/abs/2511.23257v1}.
The original author TeX, including appendices and bibliography, was
read directly. Its cyclic coefficient is checked in the programme's
complete FC1--4 proof, not copied without that calculation.
\bibitem{CH} Alain Connes, \emph{Heat Expansion and Zeta},
arXiv:2402.13082v1 (2024), Gaussian equation ft.
\url{https://arxiv.org/abs/2402.13082v1}.
The original author TeX and bibliography were read directly.
Its assumed-RH selfadjoint arithmetic interpretation is not imported.
\end{thebibliography}
\end{document}
